\documentclass[11pt,letterpaper]{article}
\usepackage{ad42010-suite}
\begin{document}
\DocFront{04}{Architecture Description Frameworks and Languages}{Specification, applicability, and conformance mapping}{An original ADF and ADL specification workbook plus a worked enterprise-software example, distinguishing description frameworks from notation and tooling.}
\InterpretationKey
\section{AD, ADF, ADL: three different artifacts}
\begin{center}
\begin{tikzpicture}[every node/.style={align=center,font=\small}, n/.style={draw=Navy,rounded corners,fill=Pale,text width=39mm,minimum height=18mm}, >=Latex]
\node[n] (a) {\textbf{ADF}\par Conventions and practices of a community};
\node[n,right=14mm of a] (b) {\textbf{ADL}\par Representations, syntax, semantics and rules};
\node[n,below=13mm of a,xshift=27mm] (c) {\textbf{AD}\par Concrete work product expressing an architecture};
\draw[->,thick] (a) -- node[left,font=\scriptsize]{organizes} (c);
\draw[->,thick] (b) -- node[right,font=\scriptsize]{supports expression} (c);
\end{tikzpicture}
\end{center}
\begin{Guidance}Do not automatically classify diagramming software, a documentation repository or a collection of templates as an ADL or ADF. Assess the claimed artifact's published scope and its constructs. An AD may use an ADF, an ADL, multiple of either, or other conventions within the standard's boundaries; merely adopting a framework does not prove AD conformance.\end{Guidance}
\section{Requirements for ADF specification}
\begin{Requirement}\FlagReq{} An ADF specification identifies the framework and scope, specifies version, and covers typical stakeholders, typical concerns and at least one viewpoint addressing those concerns. \Source{7.1.1}\end{Requirement}
\begin{longtable}{p{0.12\linewidth}p{0.54\linewidth}p{0.22\linewidth}}\toprule\textbf{Reference}&\textbf{Paraphrased ADF content}&\textbf{Evidence}\\\midrule\endhead
7.1.1(a)&Name, identity and intended scope&[URL / specification]\\
7.1.1(b)&Version consistent with local configuration rules&[version tag]\\
7.1.1(c)&At least one typical stakeholder for the applicable domain&[stakeholder catalog]\\
7.1.1(d)&At least one typical concern held by the typical stakeholders&[concern catalog]\\
7.1.1(e)&At least one viewpoint that frames the typical concerns&[viewpoint spec]\\\bottomrule\end{longtable}
\begin{Guidance}Clause 7.1.2 recommends additional framework specifications when they are expected to appear in conforming ADs: stakeholder perspectives, aspects, viewpoint-organizing formalisms, model kinds, legends, ADLs, correspondence and view methods, and version identification. Also document applicability conditions and explain consistency with the Clause 5.2 conceptual model. These recommendations should not be presented as independently mandatory 7.1.1 clauses.\end{Guidance}
\section{Reusable ADF template}
\Field{ADF name / identifier / version / steward}{[TODO]}
\Field{Domain, EoI classes and intended AD purposes}{[TODO]}
\Field{Applicability tests, triggers, exclusion conditions}{[TODO]}
\Field{Typical stakeholder classes and associated concerns}{[TODO]}
\Field{Named viewpoints and links to their 8.1 specifications}{[TODO]}
\Field{Stakeholder perspectives and aspects (recommended)}{[TODO]}
\Field{Model kinds, legends and ADLs (recommended)}{[TODO]}
\Field{View construction and correspondence methods (recommended)}{[TODO]}
\Field{Structuring formalism: domain, layer, stage, other (optional)}{[TODO]}
\Field{Map of ADF constructs to 42010 Clause 5.2 concepts (recommended)}{[TODO]}
\Field{Conditions for AD adherence / additional local rules}{[TODO]}
\section{Worked ADF: Service Architecture Description Framework}
\begin{Worked}
\textbf{Identification}: \texttt{ADF-SAD/1.0}; original example for software-intensive, network-facing business services. \textbf{Scope}: architecture decision and operational assurance reviews for production services. \textbf{Typical stakeholders}: business owner, architect, developer, operator, security and data steward. \textbf{Typical concerns}: integration boundaries, failure recovery, sensitive-data handling, runtime placement, changeability and cost.\par
\textbf{Included viewpoint specifications}: \texttt{VP-CXT/1}, \texttt{VP-FUN/1}, \texttt{VP-INFO/1}, \texttt{VP-DEP/1}, \texttt{VP-SEC/1}, \texttt{VP-OPS/1}, linked to Document 03. All are candidate resources; a particular AD selects them based on concerns.\par
\textbf{Recommended stakeholder perspectives}: business, assurance, development and operations. \textbf{Recommended aspects}: functional, behavioural, informational and structural. \textbf{Structuring formalism}: stakeholder concerns by life-cycle stage.\par
\textbf{Applicability}: security-specific viewpoints are used if the EoI processes regulated or business-sensitive data, exposes a trust boundary or has material misuse cases. Deployment view is selected if allocation affects concerns. Inapplicable viewpoints are recorded with rationale; this does \emph{not} relax universal AD requirements.\par
\textbf{Methods}: \texttt{CM-COV/1} verifies concern coverage; \texttt{CM-ALLOC/1} verifies runtime-to-deployment resolution; \texttt{VM-CXT/1} specifies how context components are produced. \textbf{ADL candidates}: UML and a documented subset of PlantUML.\par
\textbf{Local additional rule}: every approved AD records a release/environment baseline and owner; explicitly identified as an ADF-specific constraint.\par
\textbf{Source}: present suite, Document 03; not an official industry framework.
\end{Worked}
\section{Adherence to an ADF versus conformance to 42010}
\begin{Requirement}\FlagReq{} An AD's claim of conformance is assessed using the AD provisions of Clause 6. The standard separately explains how an AD can adhere to an ADF by considering the applicability of that ADF's stakeholders, perspectives, concerns, aspects, viewpoints and correspondence methods. \Source{4, 7.1}\end{Requirement}
\begin{tabularx}{\linewidth}{p{0.38\linewidth}X}\toprule\textbf{Question}&\textbf{Interpretation}\\\midrule
Does an AD use the ADF?&Catalog the selected framework, revision and relevant conditions.\\
Does it adhere to the ADF?&Map framework constructs to the AD; resolve every applicability decision; meet any additional framework rules.\\
Does the AD conform to 42010?&Independently substantiate all applicable Clause 6 mandatory provisions; framework adherence alone cannot establish this.\\
Does the ADF itself conform to 42010?&Assess its \emph{specification} against Clause 7.1, separately from the AD and the EoI.\\\bottomrule\end{tabularx}
\section{Requirements for an ADL specification}
\begin{Requirement}\FlagReq{} A specified ADL must identify the typical concerns or aspects it covers, applicable view methods, implemented model kinds, any supported viewpoints, any correspondence methods, and its version. \Source{7.2}\end{Requirement}
\begin{longtable}{p{0.12\linewidth}p{0.63\linewidth}p{0.14\linewidth}}\toprule\textbf{Ref.}&\textbf{ADL specification evidence}&\textbf{Status}\\\midrule\endhead
7.2(a)&Typical concerns and/or aspects addressed&[ ]\\
7.2(b)&At least one selected or available view method&[ ]\\
7.2(c)&One or more implemented model kinds&[ ]\\
7.2(d)&Viewpoints implemented, if any; document none where applicable&[ ]\\
7.2(e)&Correspondence methods supported, if any; document none where applicable&[ ]\\
7.2(f)&Version metadata under local rules&[ ]\\\bottomrule\end{longtable}
\subsection{Reusable ADL specification form}
\Field{ADL identity / edition / governed repository}{[TODO]}
\Field{Representation, syntax, semantics, allowed constructs}{[TODO]}
\Field{Concerns or aspects served}{[TODO]}
\Field{View methods and outputs}{[TODO]}
\Field{Model kinds implemented; supported viewpoints, if any}{[TODO]}
\Field{Correspondence methods, if any; validation guidance}{[TODO]}
\Field{Version scheme / referenced source specifications}{[TODO]}
\section{Worked ADL profile: textual UML interaction models}
\begin{Worked}\textbf{ADL example}: \texttt{ADL-INTERACTION/1.0}, a \emph{local restricted profile} for producing documented interaction views using a UML-compatible sequence notation as rendered by a text-to-diagram tool.\par
\textbf{Aspects / concerns}: behavioural interactions, timeout and responsibility flow. \textbf{Model kind}: \texttt{MK-SEQUENCE/1}. \textbf{View methods}: \texttt{VM-SEQ/1} construction, interpretation and negative-path analysis. \textbf{Implemented viewpoint}: \texttt{VP-FUN/1} only when the selection record establishes that link. \textbf{Correspondence methods}: \texttt{CM-EVENT/1} links messages to the registered interface IDs.\par
\textbf{Syntax and semantics}: participating lifelines represent externally identified runtime collaborators; solid arrows denote requests, dashed arrows responses, \texttt{alt} indicates scenario alternatives and \texttt{opt} conditional interaction. Restrict renderer features to approved constructs, and define exceptions explicitly.\par
\textbf{Version}: 1.0 for this example specification. This describes an ADL \emph{profile}, not a claim that a renderer or UML itself is certified for Clause 7.2.
\end{Worked}
\section{Using more than one ADF or ADL}
\begin{Guidance}When an AD combines framework conventions, reconcile potentially incompatible vocabularies, perspectives, aspects, model kinds, identifier schemes and correspondence rules. Maintain a crosswalk of exact concepts and denote any transformations that discard semantics. Treat profile versions and source standards as external dependencies that can change independently.\end{Guidance}
\begin{longtable}{p{0.28\linewidth}p{0.29\linewidth}p{0.30\linewidth}}\toprule\textbf{Source construct}&\textbf{Mapped AD construct}&\textbf{Equivalence / transformation / issue}\\\midrule\endhead
Framework ``actor''&Stakeholder or external entity?&Not automatically equivalent; decide by role and EoI context.\\
Diagram ``component''&View component?&Different concepts: diagram node may be an EoI constituent, while view component is a portion of the AD.\\
C4 ``view''&Architecture view?&Only if its governing viewpoint, covered concerns and requisite metadata are explicitly established.\\
UML interaction message&Interface or behavioural AD element&Map by stable identifiers using a specified correspondence method.\\\bottomrule\end{longtable}
\section{Assessment worksheet}
\Field{Artifact class for this claim: ADF or ADL}{[TODO]}
\Field{Claim scope and exact version}{[TODO]}
\Field{Clause-to-evidence mapping / unresolved gaps}{[TODO]}
\Field{Review owner, review date, decision and limitations}{[TODO]}
\LegalFooter
\end{document}
